\section{Усреднение функций}

\subsection{Основные свойства}

\subsubsection*{Ядро усреднения}

\begin{definition}
Пусть $n \in \N$, $x \in \R^n$. Положим
\[
  \omega(x) =
  \begin{cases}
    C\, e^{\frac{1}{|x|^2 - 1}}, & |x| < 1, \\[4pt]
    0, & |x| \ge 1,
  \end{cases}
  \qquad
  C = \left( \int\limits_{\ball{0}{1}} e^{\frac{1}{|x|^2 - 1}}\, dx \right)^{-1} .
\]
\end{definition}

Здесь $|x| = \sqrt{x_1^2 + \dots + x_n^2}$, а $\ball{0}{1} = \{x : |x| < 1\}$ — открытый
единичный шар. Константа $C$ подобрана так, что
\[
  \int\limits_{\R^n} \omega(x)\, dx = \int\limits_{\ball{0}{1}} \omega(x)\, dx = 1 .
\]

\begin{idea}
\textbf{Почему $C$ определена корректно.} Подынтегральная функция непрерывна, положительна
в шаре и не превосходит $e^{-1}$, а шар имеет конечный объём. Поэтому интеграл конечен
и положителен.
\end{idea}

\textbf{Свойства $\omega$:}
\begin{enumerate}
  \item $\omega \ge 0$ на $\R^n$, причём $\omega > 0$ в $\ball{0}{1}$.
  \item $\supp \omega = \overline{\ball{0}{1}}$ — компакт.
  \item $\omega$ радиальна, то есть зависит только от $|x|$. В частности,
        $\omega(-x) = \omega(x)$.
  \item $\omega \in C^\infty(\R^n)$. Это видно из записи $\omega(x) = C\, g(|x|^2)$,
        где $g(t) = e^{\frac{1}{t-1}}$ при $t < 1$ и $g(t) = 0$ при $t \ge 1$. Функция $g$
        гладкая по тем же причинам, что и шапочка из введения, а
        $|x|^2 = x_1^2+\dots+x_n^2$ — многочлен. Запись через $|x|^2$ удобна тем, что сама
        норма $|x|$ в нуле не дифференцируема.
  \item $\displaystyle\int_{\R^n} \omega\,dx = 1$.
\end{enumerate}

\begin{remark}[из тетради]
Конкретный вид $\omega$ не важен: это просто бесконечно дифференцируемая финитная функция,
и вместо неё можно взять \textbf{любую другую такую же}.
\end{remark}

\begin{idea}
В доказательствах ниже используются только свойства 1, 2, 4 и 5: гладкость,
неотрицательность, носитель в единичном шаре и интеграл, равный единице. Любое ядро
с этими свойствами подходит.

\medskip
\textbf{Что это за объект.} $\omega$ — гладкая «плотность вероятности», сосредоточенная
в единичном шаре. Свёртка функции с таким ядром заменяет значение функции в точке на
взвешенное среднее по маленькой окрестности. Отсюда название — \textbf{усреднение}.
\end{idea}

\subsubsection*{Ядро с параметром $\rho$}

\begin{definition}
Пусть $\rho > 0$. Положим
\[
  \omega_\rho(x) = \rho^{-n}\, \omega\!\left(\frac{x}{\rho}\right),
  \qquad \supp \omega_\rho = \overline{\ball{0}{\rho}} .
\]
\end{definition}

\begin{warn}
На доске и в тетради написано $\supp\omega_\rho = \ball{0}{\rho}$. Носитель — замкнутое
множество, поэтому точнее писать $\overline{\ball{0}{\rho}}$. Для рассуждений это не важно.
\end{warn}

\textbf{Интеграл сохраняется:}
\[
  \int\limits_{\ball{0}{\rho}} \omega_\rho(x)\, dx
  = \int\limits_{\ball{0}{\rho}} \rho^{-n}\, \omega\!\left(\frac{x}{\rho}\right) dx
  \;\overset{\substack{y = x/\rho\\ dx = \rho^n dy}}{=}\;
  \int\limits_{\ball{0}{1}} \rho^{-n}\, \omega(y)\, \rho^n\, dy = 1 .
\]
Замена $y = x/\rho$ переводит шар $\ball{0}{\rho}$ в $\ball{0}{1}$, а её якобиан равен $\rho^n$.

Как и $\omega$, ядро $\omega_\rho$ гладкое, неотрицательное и радиальное.

\inputfig{figures/fig-kernel-rho}

\begin{idea}
\textbf{Картинка: уменьшение $\rho$.} Аргумент $x/\rho$ сжимает шапочку до шара радиуса
$\rho$, а множитель $\rho^{-n}$ поднимает её так, чтобы интеграл остался равным $1$.
Чем меньше $\rho$, тем шапочка уже и выше: $\omega_\rho(0) = C e^{-1}\rho^{-n} \to \infty$.
При $\rho \to 0$ вся «масса» собирается в нуле, то есть $\omega_\rho$ приближает
дельта-функцию.
\end{idea}

\subsubsection*{Усреднение функции}

Пусть $\Om \subset \R^n$ — открытое ограниченное множество (в тетради: «область»)
и $u \in \Lp{1}(\Om)$.

\begin{definition}
Функция
\[
  u_\rho(x) = \int\limits_{\Om} \omega_\rho(x - y)\, u(y)\, dy
\]
называется \textbf{усреднением функции $u$ с параметром $\rho > 0$}.
\end{definition}

\begin{idea}
При фиксированном $x$ функция $y \mapsto \omega_\rho(x-y)$ — шапочка с центром в точке $x$.
Поэтому $u_\rho(x)$ — взвешенное среднее значений $u$ в $\rho$-окрестности точки $x$.
По сути, $u_\rho = \omega_\rho * u$, где $u$ продолжена нулём вне $\Om$.
\end{idea}

\paragraph{Свойства усреднения.}

\textbf{1) $u_\rho : \R^n \to \R$, то есть $u_\rho$ определена во всех точках $\R^n$.}
Интеграл конечен, поскольку $\omega_\rho$ ограничена, а $u$ суммируема:
$|\omega_\rho(x-y)\,u(y)| \le \big(\sup \omega_\rho\big)\, |u(y)|$.
Подробнее об этом рассуждении — в отступлении ниже.

\textbf{2) Интегрирование фактически идёт по $\Om \cap \ball{x}{\rho}$.}
Поскольку $\supp_y\, \omega_\rho(x-y) = \overline{\ball{x}{\rho}}$,
\[
  u_\rho(x) = \int\limits_{\Om} \omega_\rho(x - y)\, u(y)\, dy
            = \int\limits_{\Om \cap \ball{x}{\rho}} \omega_\rho(x - y)\, u(y)\, dy ,
\]
и, в частности,
\[
  u_\rho(x) = 0, \quad \text{если } \dist(x, \Om) \ge \rho .
\]

\inputfig{figures/fig-support}

\begin{idea}
\textbf{Картинка из тетради.} Вокруг $\Om$ нарисована пунктиром $\rho$-окрестность. Если шар
$\ball{x}{\rho}$ целиком лежит внутри $\Om$, то это обычное усреднение. Если $x$ около
границы, шар частично выходит наружу, и интегрирование идёт только по части шара внутри
$\Om$; это всё равно что считать $u = 0$ вне $\Om$. Если $x$ дальше $\rho$ от $\Om$, шар
с $\Om$ не пересекается, и $u_\rho(x) = 0$. Итог: \textbf{носитель $u_\rho$ лежит
в замкнутой $\rho$-окрестности $\Om$}.
\end{idea}

\textbf{3) Гладкость: $u_\rho \in C^\infty(\R^n)$.} Интеграл «наследует» гладкость ядра
$\omega_\rho$: подынтегральная функция гладкая по $x$ и суммируемая по $y$. Используется
теорема о дифференцировании интеграла Лебега по параметру из курса анализа (как сказано
на лекции, она была в конце 5-го семестра).

\inputfig{figures/fig-smoothing}

\begin{idea}
\textbf{Подробнее.} Для любого мультииндекса $\alpha$
\[
  \big|D_x^\alpha\big[\omega_\rho(x-y)\,u(y)\big]\big|
  = \big|(D^\alpha\omega_\rho)(x-y)\big|\,|u(y)|
  \le \Big(\sup_{\R^n} |D^\alpha \omega_\rho|\Big)\, |u(y)| ,
\]
и правая часть — суммируемая мажоранта, не зависящая от $x$. Значит, дифференцировать
можно под знаком интеграла:
\[
  D^\alpha u_\rho(x) = \int\limits_\Om (D^\alpha \omega_\rho)(x-y)\, u(y)\, dy ,
\]
и эти производные непрерывны. Все производные «падают» на ядро, поэтому $u$ может быть
сколь угодно негладкой.

\medskip
\textbf{Картинка из тетради.} Изображён «рваный» график $u$, под ним маленькая шапочка,
скользящая вдоль оси. В каждой точке значение $u_\rho$ — среднее $u$ под шапочкой.
Поэтому $u_\rho$ сглаживает зубцы, и получается гладкая кривая.
\end{idea}

\subsubsection*{Отступление: мажорантный признак (вопрос с лекции)}

\begin{criterion}
Пусть $f, g$ измеримы на $E$, $|f| \le g$ на $E$ и $g$ суммируема. Тогда $f$ суммируема.
\end{criterion}

\begin{corollary*}
Если $u \in \Lp{1}(E)$ и $v \in \Lp{\infty}(E)$, то $uv \in \Lp{1}(E)$.
\end{corollary*}

Действительно, $uv$ измерима и
\[
  |uv| \le |u| \cdot \norm{v}_{\Lp{\infty}(E)} \quad \text{п.\,в. на } E .
\]
Здесь $|u|$ суммируема, а $\norm{v}_{\Lp{\infty}(E)} = \esssup_E |v|$ — константа.
Значит, $uv$ суммируема.

\begin{idea}
Именно этот признак работает в свойствах 1 и 3: $\omega_\rho$ и её производные ограничены
(лежат в $\Lp{\infty}$), а $u \in \Lp{1}$.
\end{idea}

\subsubsection*{Усреднение не увеличивает норму в $\Lp{p}$}

\begin{theorem}\label{thm:norm}
Пусть $1 \le p \le \infty$, $u \in \Lp{p}(\Om)$. Тогда
\[
  \norm{u_\rho}_{\Lp{p}(\Om)} \le \norm{u}_{\Lp{p}(\Om)} .
\]
\end{theorem}

\begin{idea}
Так как $\Om$ ограничено, $\Lp{p}(\Om) \subset \Lp{1}(\Om)$ (по неравенству Гёльдера
$\norm{u}_{\Lp{1}} \le |\Om|^{1-1/p}\norm{u}_{\Lp{p}}$). Поэтому усреднение $u_\rho$
определено.

\medskip
\textbf{Смысл.} Взвешенное среднее не может быть «больше» исходной функции. Иначе говоря,
оператор $u \mapsto u_\rho$ ограничен в $\Lp{p}$ и его норма не больше $1$.
\end{idea}

\textbf{Важное утверждение (используется во всех трёх случаях).} Для любого $x \in \R^n$
\begin{equation}\label{eq:star}
  \int\limits_{\Om} \omega_\rho(x - y)\, dy \le \int\limits_{\ball{x}{\rho}} \omega_\rho(x - y)\, dy = 1 .
\end{equation}
Неравенство верно, потому что $\omega_\rho \ge 0$ и
$\int_\Om = \int_{\Om\cap \ball{x}{\rho}} \le \int_{\ball{x}{\rho}}$. Последний интеграл
после замены $z = x - y$ превращается в $\int_{\ball{0}{\rho}}\omega_\rho(z)\,dz = 1$.

Тот же результат верен, если интегрировать по $x$, а не по $y$:
$\displaystyle\int_\Om \omega_\rho(x-y)\,dx \le 1$ при любом $y$.

\begin{proof}
\textbf{1$^\circ$. Случай $p = \infty$.}
\[
  |u_\rho(x)| \le \int\limits_\Om \omega_\rho(x - y)\, |u(y)|\, dy
  \le \norm{u}_{\Lp{\infty}(\Om)} \int\limits_\Om \omega_\rho(x - y)\, dy
  \overset{\eqref{eq:star}}{\le} \norm{u}_{\Lp{\infty}(\Om)} .
\]
На втором шаге вынесли $\esssup|u|$: $|u(y)| \le \norm{u}_{\Lp{\infty}}$ при почти всех $y$.
Следовательно, $\norm{u_\rho}_{\Lp{\infty}(\Om)} \le \norm{u}_{\Lp{\infty}(\Om)}$.

\begin{warn}
На доске замена записана как $z = y - x$, а в тетради — как $z = x - y$. Удобнее второй
вариант: при нём сразу $\omega_\rho(x-y) = \omega_\rho(z)$. При $z = y - x$ получается
$\omega_\rho(-z)$, и нужно дополнительно сослаться на чётность ядра. Также в тетради
в последнем интеграле написано $\omega(z)$ вместо $\omega_\rho(z)$ — это описка.
\end{warn}

\textbf{2$^\circ$. Случай $p = 1$.}
\begin{multline*}
  \int\limits_\Om |u_\rho(x)|\, dx
  \le \int\limits_\Om dx \int\limits_\Om \omega_\rho(x - y)\, |u(y)|\, dy
  \overset{\text{Фубини--Тонелли}}{=} \\
  = \int\limits_\Om |u(y)|\, dy
    \underbrace{\int\limits_\Om \omega_\rho(x - y)\, dx}_{\le 1 \text{ по } \eqref{eq:star}}
  \le \int\limits_\Om |u(y)|\, dy .
\end{multline*}
Значит, $\norm{u_\rho}_{\Lp{1}(\Om)} \le \norm{u}_{\Lp{1}(\Om)}$.

\begin{idea}
Менять порядок интегрирования можно по \textbf{теореме Тонелли}: подынтегральная функция
$\omega_\rho(x-y)|u(y)|$ неотрицательна, а для неотрицательных функций перестановка
допустима без предварительной проверки суммируемости.
\end{idea}

\textbf{3$^\circ$. Случай $1 < p < \infty$.} Пусть $q$ — сопряжённый показатель:
$\frac1p + \frac1q = 1$.

\begin{reminder}[неравенство Гёльдера]
Если $u \in \Lp{p}(\Om)$ и $v \in \Lp{q}(\Om)$, то $uv \in \Lp{1}(\Om)$ и
\[
  \int\limits_\Om |uv|\, dx \le \left(\int\limits_\Om |u|^p dx\right)^{1/p}
                                \left(\int\limits_\Om |v|^q dx\right)^{1/q} .
\]
\end{reminder}

Основной приём — разложить ядро на два множителя:
$\omega_\rho = \omega_\rho^{1/q}\cdot\omega_\rho^{1/p}$, поскольку $\frac1q + \frac1p = 1$.
Тогда по Гёльдеру
\begin{multline*}
  |u_\rho(x)| \le \int\limits_\Om \omega_\rho^{\frac1q}(x - y)\cdot \omega_\rho^{\frac1p}(x-y)\,|u(y)|\, dy \\
  \le \left( \int\limits_\Om \omega_\rho(x - y)\, dy \right)^{\frac1q}
      \left( \int\limits_\Om \omega_\rho(x - y)\, |u(y)|^p\, dy \right)^{\frac1p} .
\end{multline*}
Первая скобка не больше $1$ по \eqref{eq:star}. Возводя в степень $p$, получаем
\[
  |u_\rho(x)|^p \le \int\limits_\Om \omega_\rho(x - y)\, |u(y)|^p\, dy .
\]
Проинтегрируем по $x$ и снова применим Фубини--Тонелли и \eqref{eq:star}:
\begin{multline*}
  \norm{u_\rho}^p_{\Lp{p}(\Om)} = \int\limits_\Om |u_\rho(x)|^p dx
  \le \int\limits_\Om dx \int\limits_\Om \omega_\rho(x-y)\,|u(y)|^p\,dy \\
  = \int\limits_\Om |u(y)|^p\,dy \underbrace{\int\limits_\Om \omega_\rho(x-y)\,dx}_{\le 1}
  \le \norm{u}^p_{\Lp{p}(\Om)} . \qquad\qedhere
\end{multline*}
\end{proof}

\begin{idea}
\textbf{Идея случая 3$^\circ$.} Он сводится к случаю $p = 1$, применённому к функции
$|u|^p$. Разложение $\omega_\rho = \omega_\rho^{1/q}\omega_\rho^{1/p}$ нужно, чтобы
«забрать» часть веса в каждый сомножитель Гёльдера. Фактически это неравенство Йенсена
для выпуклой функции $t \mapsto |t|^p$ и вероятностной меры $\omega_\rho(x-y)\,dy$.
\end{idea}

\subsubsection*{Утверждение: равномерная сходимость на компактах}

\begin{proposition}\label{prop:unif}
Пусть $u \in C(\clOm)$, $K$ — компакт и $K \subset \Om$. Тогда
\[
  u_\rho \unif u \ \text{ на } K \quad \text{при } \rho \to 0 ,
\]
то есть $u_{\rho_k} \unif u$ на $K$ для любой последовательности $\rho_k > 0$, $\rho_k \to 0$.
\end{proposition}

\inputfig{figures/fig-compact}

\begin{proof}
\emph{Шаг 1: компакт отделён от границы.} Граница
$\dOm = \clOm \setminus \Om = \clOm \cap (\complement\Om)$ — всегда замкнутое множество.
Поскольку $\Om$ открыто, $\dOm \cap \Om = \varnothing$, а значит, $K \cap \dOm = \varnothing$.
Расстояние между непересекающимися компактом и замкнутым множеством положительно:
\[
  d = \dist(K, \dOm) > 0 .
\]

\emph{Шаг 2: при малых $\rho$ шары не вылезают из $\Om$.} Пусть $\rho < d$ и $x \in K$.
Тогда $\ball{x}{\rho} \subset \Om$, и поэтому
\[
  \int\limits_\Om \omega_\rho(x-y)\,dy
  = \int\limits_{\Om \cap \ball{x}{\rho}} \omega_\rho(x-y)\,dy
  = \int\limits_{\ball{x}{\rho}} \omega_\rho(x-y)\,dy = 1 .
\]

\begin{idea}
Почему $\ball{x}{\rho} \subset \Om$: если бы какая-то точка $z \in \ball{x}{\rho}$ не лежала
в $\Om$, то на отрезке $[x, z]$ нашлась бы точка границы $\dOm$ на расстоянии меньше
$\rho < d$ от $x \in K$. Это противоречит определению $d$. Обратите внимание: здесь
неравенство \eqref{eq:star} превращается в равенство. Около компакта усреднение берётся
по полному шару.
\end{idea}

\emph{Шаг 3: оценка разности.} Пусть $x \in K$, $\rho < d$. Представим $u(x)$ как
$u(x) \cdot 1$, заменив единицу интегралом из шага 2:
\[
  u(x) - u_\rho(x)
  = u(x)\int\limits_\Om \omega_\rho(x-y)\,dy - \int\limits_\Om \omega_\rho(x-y)\,u(y)\,dy
  = \int\limits_\Om \omega_\rho(x-y)\,\big(u(x) - u(y)\big)\,dy ,
\]
откуда
\[
  |u(x) - u_\rho(x)| \le \int\limits_{\ball{x}{\rho}} \omega_\rho(x-y)\,|u(x) - u(y)|\,dy .
\]

\emph{Шаг 4: равномерная непрерывность.} Возьмём $\eps > 0$. Найдётся $\delta > 0$ такое,
что $|u(x) - u(y)| < \eps$ при $|x - y| < \delta$. При $\rho < \min(d, \delta)$ все
$y \in \ball{x}{\rho}$ удовлетворяют $|x-y| < \delta$, поэтому
\[
  |u(x) - u_\rho(x)| \le \eps \int\limits_{\ball{x}{\rho}} \omega_\rho(x-y)\,dy = \eps
  \qquad \text{для всех } x \in K .
\]
Оценка не зависит от $x \in K$, значит, сходимость равномерная.
\end{proof}

\begin{warn}
В тетради написано «$u$ равномерно непрерывна на $K$». Этого недостаточно: точка $y$
пробегает шар $\ball{x}{\rho}$ и может не лежать в $K$. Правильно так: $\Om$ ограничено,
поэтому $\clOm$ — компакт, и $u \in C(\clOm)$ равномерно непрерывна на всём $\clOm$
(теорема Кантора).
\end{warn}

\begin{idea}
\textbf{Идея доказательства.} Вдали от границы усреднение — честное среднее по шару
с полным весом $1$. Поэтому $u - u_\rho$ — среднее отклонений $u(x)-u(y)$ по маленькому
шару, а по непрерывности эти отклонения малы.
\end{idea}

\begin{corollary*}
Если $u \in C(\clOm)$, то $u_\rho(x) \to u(x)$ при $\rho \to 0$ для каждого $x \in \Om$.
\end{corollary*}

\begin{proof}
Точка — это компакт: достаточно взять $K = \{x\}$.
\end{proof}

\subsubsection*{Сходимость усреднений в $\Lp{p}$}

\begin{theorem}\label{thm:conv}
Пусть $u \in \Lp{p}(\Om)$, $1 \le p < \infty$. Тогда
\[
  u_\rho \to u \ \text{ в } \Lp{p}(\Om) \quad \text{при } \rho \to 0 .
\]
\end{theorem}

\begin{idea}
\textbf{Почему $p < \infty$.} При $p = \infty$ теорема неверна. Все $u_\rho$ непрерывны,
а равномерный предел непрерывных функций непрерывен, поэтому разрывную функцию
(например, ступеньку из примера ниже) нельзя приблизить усреднениями в $\Lp{\infty}$.
\end{idea}

\begin{proof}
Сначала докажем для непрерывных функций, затем перенесём на все $\Lp{p}$ по плотности.

\textbf{1) Случай $u \in C(\clOm)$.} Возьмём компакт $K \subset \Om$ и разобьём
$\Om = K \cup (\Om \setminus K)$:
\[
  \norm{u - u_\rho}^p_{\Lp{p}(\Om)}
  = \int\limits_K |u - u_\rho|^p\,dx + \int\limits_{\Om\setminus K} |u - u_\rho|^p\,dx .
\]
Обозначим $M = \norm{u}_{C(\clOm)} = \norm{u}_{\Lp{\infty}(\Om)}$. По теореме~\ref{thm:norm}
(случай $p=\infty$) $\norm{u_\rho}_{\Lp{\infty}(\Om)} \le M$, поэтому $|u - u_\rho| \le 2M$
на $\Om$. Отсюда
\[
  \norm{u - u_\rho}^p_{\Lp{p}(\Om)}
  \le \sup_K |u - u_\rho|^p \cdot \mu(K) + (2M)^p \cdot \mu(\Om \setminus K) .
\]
Дальше по данному $\eps > 0$:
\begin{itemize}
  \item \textbf{сначала выбираем $K$} так, чтобы $(2M)^p\,\mu(\Om\setminus K) < \eps/2$.
        Это всегда возможно;
  \item \textbf{затем выбираем $\rho$}: так как $u_\rho \unif u$ на $K$
        (утверждение~\ref{prop:unif}), при всех достаточно малых $\rho$ выполнено
        $\sup_K |u - u_\rho|^p\,\mu(K) < \eps/2$.
\end{itemize}
Итого $\norm{u - u_\rho}^p_{\Lp{p}} < \eps$ для всех достаточно малых $\rho$.

\begin{idea}
\textbf{Почему $K$ можно выбрать.} Возьмём
$K_m = \{x \in \Om : \dist(x, \dOm) \ge \tfrac1m\}$. Это компакты (замкнуты и ограничены),
$K_m \subset \Om$, они возрастают и в объединении дают $\Om$. Поскольку $\mu(\Om) < \infty$,
по непрерывности меры $\mu(\Om \setminus K_m) \to 0$.

\medskip
\textbf{Порядок выбора важен:} $K$ выбирается раньше $\rho$ и от $\rho$ не зависит. Первое
слагаемое контролируется равномерной сходимостью на фиксированном компакте, второе —
малостью «каёмки» $\Om \setminus K$ около границы. Около границы сходимости может и не
быть, но там мало меры.
\end{idea}

\textbf{2) Общий случай $u \in \Lp{p}(\Om)$.} Пусть $\eps > 0$. Непрерывные функции плотны
в $\Lp{p}(\Om)$ при $1 \le p < \infty$, поэтому найдётся $g \in C(\clOm)$
с $\norm{u - g}_{\Lp{p}} < \eps/4$. Вставим $g$ и $g_\rho$:
\[
  \norm{u - u_\rho}_{\Lp{p}} = \norm{(u - g) + (g - g_\rho) + (g_\rho - u_\rho)}_{\Lp{p}}
  \le \norm{u - g}_{\Lp{p}} + \norm{g - g_\rho}_{\Lp{p}} + \norm{g_\rho - u_\rho}_{\Lp{p}} .
\]
Усреднение линейно по $u$, поэтому $g_\rho - u_\rho = (g - u)_\rho$. По теореме~\ref{thm:norm}
$\norm{(g-u)_\rho}_{\Lp{p}} \le \norm{g - u}_{\Lp{p}}$. Значит,
\[
  \norm{u - u_\rho}_{\Lp{p}}
  \le 2\underbrace{\norm{u - g}_{\Lp{p}}}_{< \eps/4}
   + \underbrace{\norm{g - g_\rho}_{\Lp{p}}}_{< \eps/2 \text{ при малых } \rho} < \eps .
\]
Второе слагаемое мало при малых $\rho$ по пункту 1, применённому к $g$.
\end{proof}

\begin{idea}
\textbf{Схема «$\eps/3$».} Это стандартный приём: (а) доказать сходимость на плотном классе
хороших функций; (б) иметь \textbf{равномерную по $\rho$} оценку нормы операторов, здесь —
теорему~\ref{thm:norm}. Тогда сходимость переносится на всё пространство. Теорема~\ref{thm:norm}
нужна именно здесь: без неё слагаемое $\norm{g_\rho - u_\rho}$ нечем было бы контролировать.

\medskip
Плотность: из курса теории меры известно, что в $\Lp{p}(\Om)$, $1 \le p < \infty$, плотны
даже непрерывные функции с компактным носителем в $\Om$. Продолженные нулём, они лежат
в $C(\clOm)$.
\end{idea}

\subsubsection*{Пример: усреднение ступеньки}

\begin{example}
Пусть $n = 1$, $\Om = (-1, 1)$,
\[
  u(x) = \begin{cases} 0, & x < 0, \\ 1, & x \ge 0, \end{cases}
  \qquad \rho \in \left\{1, \tfrac12, \tfrac14, \tfrac1{10}\right\}.
\]
Ядро: $\omega(x) = C e^{\frac{1}{x^2-1}}$ при $|x| < 1$ и $\omega(x) = 0$ при $|x| \ge 1$;
$\omega_\rho(x) = \frac1\rho\,\omega\!\left(\frac x\rho\right)$.
\end{example}

\begin{warn}
В тетради у $\omega$ пропущена нормирующая константа $C$. Для $n = 1$ получается
$C \approx 2{,}25$, так как $\int_{-1}^1 e^{\frac{1}{x^2-1}}dx \approx 0{,}444$. Например,
высота $\omega_{1/2}$ в нуле равна $2Ce^{-1} \approx 1{,}66$, а носитель — $[-\tfrac12, \tfrac12]$.
\end{warn}

\inputfig{figures/fig-step-setup}

Поскольку $u = 0$ на $(-1,0)$,
\[
  u_\rho(x) = \int\limits_{-1}^{1} \omega_\rho(x - y)\,u(y)\,dy = \int\limits_0^1 \omega_\rho(x - y)\,dy .
\]

\begin{idea}
\textbf{Как это посчитать.} После замены $z = x - y$ получаем
$u_\rho(x) = \int_{x-1}^{x} \omega_\rho(z)\,dz = \Phi_\rho(x) - \Phi_\rho(x-1)$, где
$\Phi_\rho(t) = \int_{-\infty}^t \omega_\rho$ — «гладкая ступенька»: $\Phi_\rho = 0$
при $t \le -\rho$, $\Phi_\rho = 1$ при $t \ge \rho$, а на $[-\rho, \rho]$ она гладко растёт.
Поэтому при $\rho \le \tfrac12$:

\medskip
\begin{center}
\begin{tabular}{ll}
\toprule
участок & значение $u_\rho(x)$ \\
\midrule
$x \le -\rho$ & $0$ \\
$-\rho \le x \le \rho$ & гладко растёт от $0$ до $1$ (в точке $x=0$ равно $\tfrac12$) \\
$\rho \le x \le 1-\rho$ & $1$ \\
$1-\rho \le x \le 1+\rho$ & гладко убывает от $1$ до $0$ (в точке $x = 1$ равно $\tfrac12$) \\
$x \ge 1+\rho$ & $0$ \\
\bottomrule
\end{tabular}
\end{center}
\end{idea}

\inputfig{figures/fig-step-mollified}

\paragraph{Графики из тетради.}
\begin{itemize}
  \item \textbf{Сверху:} $u_\rho$ около скачка в нуле. Вместо разрыва — гладкий подъём
        на отрезке $[-\rho, \rho]$, и чем меньше $\rho$, тем он круче.
  \item \textbf{Снизу:} $u_\rho$ на всей прямой.
        \begin{itemize}
          \item При $\rho = \tfrac12$ — «горб»: подъём на $[-\tfrac12, \tfrac12]$, значение
                $1$ ровно в точке $\tfrac12$, затем спуск до нуля к $x = \tfrac32$.
          \item При $\rho = \tfrac14$ — «трапеция»: подъём на $[-\tfrac14, \tfrac14]$,
                плато $u_\rho = 1$ на $[\tfrac14, \tfrac34]$, спуск на $[\tfrac34, \tfrac54]$.
        \end{itemize}
  \item При $\rho = 1$ функция $u_\rho$ вообще не достигает единицы: окно усреднения шире,
        чем участок, где $u = 1$.
\end{itemize}

\begin{idea}
\textbf{Что здесь видно.}
\begin{itemize}
  \item \textbf{Гладкость (свойство 3):} разрывная ступенька превратилась в бесконечно
        гладкую функцию.
  \item \textbf{Носитель (свойство 2):} $u_\rho$ «вылезает» за пределы $\Om$ на расстояние
        $\rho$ (правый спуск заканчивается в $1 + \rho$). Спуск около $x = 1$ возникает
        потому, что вне $\Om$ функция $u$ считается нулём.
  \item \textbf{Сходимость в $\Lp{p}$ (теорема~\ref{thm:conv}):} $u_\rho$ отличается от $u$
        только на отрезках длины порядка $\rho$ около $0$ и около $1$, и там
        $|u - u_\rho| \le 1$. Поэтому $\norm{u - u_\rho}^p_{\Lp{p}(\Om)} \lesssim \rho \to 0$.
  \item \textbf{Нет равномерной сходимости:} $u_\rho(0) = \tfrac12$ при всех малых $\rho$,
        а $u(0) = 1$. Утверждение о поточечной и равномерной сходимости требует
        непрерывности $u$, и ступенька показывает, что без неё оно неверно.
\end{itemize}
\end{idea}

\lecture{Лекция 2}{15.09.2026}

\subsubsection*{Следствия: плотность гладких функций в $\Lp{p}$}

\begin{corollary}\label{cor:dense-smooth}
$C^\infty(\clOm)$ плотно в $\Lp{p}(\Om)$, $1 \le p < \infty$. В качестве приближений можно
брать усреднения $u_\rho$.
\end{corollary}

\begin{proof}
По свойству 3 $u_\rho \in C^\infty(\R^n)$, поэтому сужение $u_\rho$ на $\Om$ лежит
в $C^\infty(\clOm)$: все производные непрерывно продолжаются до границы.
По теореме~\ref{thm:conv} $u_\rho \to u$ в $\Lp{p}(\Om)$.
\end{proof}

\begin{warn}
В тетради написано «$1 \le p = \infty$». Должно быть $1 \le p < \infty$: при $p = \infty$
утверждение неверно (см. пример со ступенькой).
\end{warn}

\begin{corollary}\label{cor:dense-c0}
$C_0^\infty(\Om)$ плотно в $\Lp{p}(\Om)$, $1 \le p < \infty$.
\end{corollary}

\begin{reminder}[что такое $C_0^\infty(\Om)$]
Если $u \in C_0^\infty(\Om)$, то, продолженная нулём, $u \in C^\infty(\R^n)$, и
\[
  \supp u \Subset \Om \iff \supp u \text{ — компакт и } \supp u \subset \Om .
\]
Знак $\Subset$ означает «строго внутри»: множество вместе с замыканием лежит в $\Om$.
\end{reminder}

\begin{idea}
\textbf{Разница со следствием~\ref{cor:dense-smooth}.} Функции из $C^\infty(\clOm)$ могут
быть ненулевыми вплоть до границы. Функции из $C_0^\infty(\Om)$ обязаны обращаться в ноль
в некоторой полосе около границы. Обычное усреднение $u_\rho$ этого не гарантирует: его
носитель, наоборот, вылезает за $\Om$. Поэтому функцию нужно сначала «обрезать» внутри $\Om$.
\end{idea}

\begin{proof}
Пусть $u \in \Lp{p}(\Om)$ и $\eps > 0$.

\emph{Шаг 1: обрезка.} Для компакта $K \subset \Om$ положим
\[
  u_K = u\,\chi_K = \begin{cases} u, & \text{на } K, \\ 0, & \text{вне } K. \end{cases}
\]
Тогда
\[
  \norm{u - u_K}_{\Lp{p}(\Om)}
  = \left(\int\limits_\Om |u - u_K|^p dx\right)^{1/p}
  = \left(\int\limits_{\Om\setminus K} |u - u_K|^p dx\right)^{1/p}
  = \left(\int\limits_{\Om\setminus K} |u|^p dx\right)^{1/p} .
\]
Поскольку $|u|^p \in \Lp{1}(\Om)$, по \textbf{абсолютной непрерывности интеграла Лебега}
(как функции множества) существует $\delta > 0$ такое, что
\[
  \int\limits_{E_0} |u|^p\,dx < \left(\frac\eps2\right)^p ,
  \quad \text{как только } \mu(E_0) < \delta .
\]
Выберем компакт $K \subset \Om$ с $\mu(\Om \setminus K) < \delta$; это возможно,
см. доказательство теоремы~\ref{thm:conv}. Тогда $\norm{u - u_K}_{\Lp{p}} < \eps/2$.

\emph{Шаг 2: усреднение обрезанной функции.} Рассмотрим
\[
  (u_K)_\rho(x) = \int\limits_\Om \omega_\rho(x-y)\,u_K(y)\,dy
                = \int\limits_{K \cap \ball{x}{\rho}} \omega_\rho(x-y)\,u_K(y)\,dy ,
\]
поскольку $u_K = 0$ вне $K$, а $\supp_y \omega_\rho(x-y) = \overline{\ball{x}{\rho}}$.
Если $\rho < \dist(K, \dOm)$, то
\[
  (u_K)_\rho(x) = 0 \quad \text{для всех } x \text{ с } \dist(x, K) \ge \rho ,
\]
то есть носитель $(u_K)_\rho$ лежит в замкнутой $\rho$-окрестности компакта $K$ — она сама
компакт и целиком содержится в $\Om$ (картинка в тетради: $K$ внутри $\Om$, вокруг $K$
пунктиром $\rho$-окрестность, ср. рис.~\ref{fig:compact}). Вместе с гладкостью (свойство 3)
это даёт
\[
  (u_K)_\rho \in C_0^\infty(\Om), \qquad \rho < \dist(K, \dOm) \quad\text{— финитная функция.}
\]

\emph{Шаг 3: оценка.} По теореме~\ref{thm:conv} $(u_K)_\rho \to u_K$ в $\Lp{p}(\Om)$,
поэтому найдётся $\rho_0 > 0$ такое, что $\norm{u_K - (u_K)_\rho}_{\Lp{p}} < \eps/2$
при $\rho < \rho_0$. Тогда
\[
  \norm{u - (u_K)_\rho}_{\Lp{p}}
  \le \norm{u - u_K}_{\Lp{p}} + \norm{u_K - (u_K)_\rho}_{\Lp{p}}
  < \frac\eps2 + \frac\eps2 = \eps
\]
для $\rho < \min\big(\rho_0,\ \dist(K, \dOm)\big)$.
\end{proof}

\begin{idea}
\textbf{Идея.} Сначала отрезаем от $u$ «каёмку» у границы — это стоит мало в $\Lp{p}$
из-за абсолютной непрерывности интеграла. Затем сглаживаем оставшуюся функцию усреднением
с радиусом меньше расстояния до границы. Сглаживание размывает носитель всего на $\rho$,
и он остаётся строго внутри $\Om$.
\end{idea}

\begin{remark}
Всё это — для $p < \infty$. В $\Lp{\infty}$ не плотны даже непрерывные функции.
\end{remark}

\subsubsection*{Локально суммируемые функции}

\begin{definition}
$u \in L^p_{\mathrm{loc}}(\Om)$ («$u$ локально принадлежит $\Lp{p}$»), если
\[
  u \in \Lp{p}(\Om') \quad \forall\, \Om' \Subset \Om .
\]
Случай $p = 1$ даёт класс $\Lloc(\Om)$ локально суммируемых функций — именно он нужен
дальше в курсе.
\end{definition}

\begin{example}
$\Om = \R$: функции, суммируемые на отрезках, лежат в $\Lloc{}(\R)$.
В частности, константа $\in \Lloc{}(\R)$, хотя она не суммируема на всей прямой.
\end{example}

\begin{example}
$\Om = (0,1)$: $\dfrac1x \in \Lloc{}\big((0,1)\big)$, хотя
$\int_0^1 \frac{dx}{x} = +\infty$ и $\tfrac1x \notin \Lp{1}((0,1))$.
\end{example}

\inputfig{figures/fig-loc-1x}

\begin{remark}[иерархия классов]
\[
  C_0^\infty(\Om) \text{ — самые хорошие функции (но самый узкий класс)},
\]
\[
  \Lloc{}(\Om) \text{ — самые плохие функции (но самый широкий класс)} .
\]
\end{remark}

\subsubsection*{Основная лемма вариационного исчисления (Дюбуа--Реймона)}

\begin{lemma}[Дюбуа--Реймон]\label{lem:dbr}
Пусть $\Om$ — открытое множество (не обязательно ограниченное),
\[
  f \in \Lloc{}(\Om), \qquad
  \int\limits_\Om f\,u\,dx = 0 \quad \forall u \in C_0^\infty(\Om) .
\]
Тогда $f = 0$ почти всюду на $\Om$.
\end{lemma}

\begin{idea}
Пометка с лекции: здесь основная борьба идёт за \emph{классы функций}. Пробные функции $u$
берутся из самого узкого класса $C_0^\infty$, а вывод делается про $f$ из самого широкого
класса $\Lloc{}$. Чем уже класс пробных функций, тем сильнее утверждение.
\end{idea}

\begin{proof}
\emph{Шаг 1: локализация.} Пусть $K \subset \Om$ — компакт. Найдётся открытое ограниченное
$\widetilde\Om$ с
\[
  K \subset \widetilde\Om \Subset \Om .
\]
Так как $f \in \Lloc{}(\Om)$, то $f \in \Lp{1}(\widetilde\Om)$.

\inputfig{figures/fig-nested}

\emph{Шаг 2: пробная функция.} Положим
\[
  g = \begin{cases} \operatorname{sign} f, & \text{на } K, \\ 0, & \text{вне } K, \end{cases}
  \qquad\text{то есть } g = \operatorname{sign} f \cdot \chi_K .
\]
Это произведение измеримых функций, поэтому $g$ измерима в $\Om$. Её усреднение
\[
  g_\rho(x) = \int\limits_{\widetilde\Om} \omega_\rho(x-y)\,g(y)\,dy
            = \int\limits_{K} \omega_\rho(x-y)\,g(y)\,dy ,
\]
и, как в следствии~\ref{cor:dense-c0},
\[
  g_\rho \in C_0^\infty(\widetilde\Om) \subset C_0^\infty(\Om),
  \qquad \text{если } \rho < \dist(K, \partial\widetilde\Om) .
\]

\emph{Шаг 3: сходимость почти всюду.} По теореме~\ref{thm:conv} $g_\rho \to g$
в $\Lp{1}(\widetilde\Om)$, значит, $g_\rho \to g$ по мере в $\widetilde\Om$. По теореме
Рисса существует последовательность $\rho_k \to 0$ с
\[
  g_{\rho_k} \to g \quad \text{п.\,в. в } \widetilde\Om
  \qquad\Longrightarrow\qquad
  f\,g_{\rho_k} \to f\,g \quad \text{п.\,в. в } \widetilde\Om .
\]

\emph{Шаг 4: мажоранта (нужная ветвь рассуждений).} По построению $g \in \Lp{\infty}(\widetilde\Om)$
и $|g| \le 1$ в $\widetilde\Om$. По теореме~\ref{thm:norm} (случай $p=\infty$)
\[
  \norm{g_\rho}_{\Lp{\infty}} \le \norm{g}_{\Lp{\infty}} \le 1
  \ \Longrightarrow\ |g_\rho| \le 1 \ \text{п.\,в. в } \widetilde\Om
  \ \Longrightarrow\ |g_\rho|\cdot|f| \le |f| \ \text{п.\,в. на } \widetilde\Om ,
\]
а $f \in \Lloc{}(\Om)$ даёт $|f| \in \Lp{1}(\widetilde\Om)$. Итого по теореме
Лебега о мажорированной сходимости
\[
  \int\limits_{\widetilde\Om} f\,g_{\rho_k}\,dx \longrightarrow \int\limits_{\widetilde\Om} f\,g\,dx .
\]

\emph{Шаг 5: предел равен нулю.} Но $g_{\rho_k} \in C_0^\infty(\Om)$, поэтому по условию леммы
\[
  \int\limits_{\widetilde\Om} f\,g_{\rho_k}\,dx = \int\limits_{\Om} f\,g_{\rho_k}\,dx = 0 .
\]
Следовательно,
\[
  0 = \int\limits_{\widetilde\Om} f\,g\,dx = \int\limits_K f\,\operatorname{sign} f\,dx
    = \int\limits_K |f|\,dx
  \qquad\Longrightarrow\qquad f = 0 \ \text{п.\,в. на } K .
\]

\emph{Шаг 6: от компакта ко всей области.} По теореме о представлении открытого множества
ячейками $\Om = \bigcup_{k=1}^{\infty} \overline{P_k}$, где $P_k$ — замкнутые ячейки, то есть
компакты. Тогда
\[
  \Om[f \neq 0] = \bigcup_{i \in \N} \overline{P_i}[f \neq 0]
  \ \Longrightarrow\
  \mu\big(\Om[f \neq 0]\big) \le \sum_{i=1}^{\infty} \mu\big(\overline{P_i}[f \neq 0]\big) = 0 ,
\]
поскольку каждое слагаемое нулевое по шагу 5. Значит, $f = 0$ п.\,в. на $\Om$.
\end{proof}

\begin{idea}
\textbf{Зачем $\operatorname{sign} f$.} Пробная функция подбирается так, чтобы
$f\,g = |f|$ на $K$: тогда равенство нулю интеграла сразу даёт $f = 0$. Но
$\operatorname{sign} f \cdot \chi_K$ разрывна и не годится как пробная функция, поэтому её
сначала сглаживают усреднением — отсюда переход к $g_\rho$ и предельный переход по Лебегу.
\end{idea}
